\section{Síntesis y ampliación}

\subsection{De token relation a datapoint relation}

La visión unificada que más vale conservar de esta clase es la siguiente: la attention no pertenece por naturaleza a los token del lenguaje; es una plantilla de cómputo que «permite a un objeto actualizar su representación a partir de un conjunto de otros objetos». El Transformer aprende relaciones entre tokens; NPT convierte cada datapoint en un objeto y, mediante reshape, razona alternando entre dos ejes: datapoints y attributes.

\begin{importantbox}{Doce conclusiones centrales}
\begin{enumerate}
  \item La self-attention hace que query, key y value provengan del mismo conjunto, y aprende las relaciones internas de la entrada.
  \item La multi-head attention ofrece varios subespacios de proyección, pero no garantiza que cada head pueda recibir un nombre interpretable por una persona.
  \item Teacher forcing usa gold prefixes para entrenar en paralelo todas las next-token positions.
  \item La causal mask impide que el modelo lea los targets futuros; es el límite que asegura la corrección de la likelihood.
  \item NPT sigue siendo una parametric neural network; «non-parametric» significa que la inference depende explícitamente de los datapoints de entrenamiento.
  \item $(X,M)$ coloca la observed evidence y las masked queries en una misma dataset representation.
  \item ABD aplica attention entre rows; ABA aplica attention entre los attributes de cada fila.
  \item La row permutation equivariance garantiza que el orden de los datapoints no aporte semántica espuria.
  \item El feature masking proporciona regularization; el target masking entrena el relational lookup.
  \item La evidencia de UCI average-rank respalda que el modelo es competitivo, pero no que supere a los demás en todas las tareas.
  \item Las intervenciones de corruption y duplicate demuestran que, en varias tareas, el mecanismo entre datapoints realmente se utiliza.
  \item La quadratic datapoint attention es la limitación central; a futuro se requieren retrieval, sparsity o representative points.
\end{enumerate}
\end{importantbox}

\subsection{Un marco de decisión para la ingeniería}

Para decidir si conviene adoptar un mecanismo al estilo NPT en una tarea nueva, pueden plantearse cinco preguntas en orden: primero, si los otros examples realmente aportan información adicional durante la inference; segundo, si es posible acceder a esos examples de forma legítima, segura y con baja latencia; tercero, si el lote o el retrieval set cubre los neighbors relevantes; cuarto, si un baseline común de tree, k-NN o retrieval ya es suficiente; quinto, si las pruebas de corruption e intervention pueden demostrar que el modelo no depende únicamente de sus propias features o de un shortcut.

\begin{knowledgebox}{Lista de verificación práctica}
Primero, implementar un baseline per-row sólido; después, incorporar explicit retrieval; luego, comparar k-NN fijo, learned retriever y full datapoint attention; por último, realizar row corruption, neighbor removal, label intervention, análisis de sensibilidad a la composición del lote y OOD evaluation. Solo cuando mejoren a la vez el desempeño y la evidencia sobre el mecanismo debe atribuirse la ganancia al razonamiento entre datapoints.
\end{knowledgebox}

\subsection{Lecturas adicionales}

Se recomienda leer en el orden «mecanismo--arquitectura--validación»: primero, volver a \emph{Attention Is All You Need} para comprender el causal training; después, leer la sección 2 y los apéndices del artículo de NPT de Kossen et al. para dominar los tensor shapes y el masking; a continuación, contrastar Deep Sets, Set Transformer, Neural Relational Inference y Neural Processes, comparando permutation symmetry, message passing, context conditioning y uncertainty assumptions. En la práctica de ingeniería conviene seguir de cerca la sparse attention y el retrieval-augmented modeling, porque atacan directamente el quadratic bottleneck de NPT sobre el conjunto de datos completo.

\teachervoice{El cierre del ponente no afirma que el problema esté resuelto; deja el scaling y las tareas más ricas como direcciones abiertas. La síntesis más prudente es esta: NPT demuestra que «entrenar un algoritmo de predicción que lee el conjunto de datos» funciona y aporta experimentos poco comunes a nivel de mecanismo; al mismo tiempo, convierte la cobertura de la recuperación, la complejidad computacional y la fuga de datos en problemas de sistema ineludibles.}

\end{document}
